\documentclass[12pt, letterpaper]{article}
\usepackage[margin=1.2in]{geometry}
\usepackage{ebgaramond}
\usepackage[T1]{fontenc}
\usepackage{microtype}
\usepackage{setspace}
\usepackage{parskip}
\usepackage{titling}
\usepackage{fancyhdr}

\setlength{\parindent}{2em}
\onehalfspacing

\pagestyle{fancy}
\fancyhf{}
\fancyhead[L]{\small\textit{The Tape}}
\fancyhead[R]{\small\textit{NASA Mission Series v1.9}}
\fancyfoot[C]{\small\thepage}
\renewcommand{\headrulewidth}{0.4pt}

\title{\Huge\textbf{The Tape}\\[0.5em]\large\textit{A short story from Mission 1.9 --- Explorer 3}}
\author{NASA Mission Series\\[0.3em]\small Miles Tiberius Foxglove --- Mukilteo, WA}
\date{March 30, 2026}

\begin{document}
\maketitle
\thispagestyle{empty}

\vspace{1em}
\noindent\rule{\textwidth}{0.4pt}
\vspace{1em}

\noindent I am a loop of magnetic oxide on a strip of Mylar, wound tight inside a metal box, bolted to the instrument shelf of a satellite that weighs as much as a small dog. I am four meters long. I am six millimeters wide. I was built in a basement laboratory at the University of Iowa by a graduate student named George Ludwig, who wound me by hand and tested me in a vacuum chamber and prayed over me in the language of engineers, which is repeated measurement. I am the smallest memory in space.

I do not know that I am in space. I know only that the capstan turns and my surface moves past the head, and as it moves, the head's magnetic field rearranges the oxide particles on my coating into patterns. Each pattern is a number. Each number is a count. Each count is a particle --- a proton or electron that struck the Geiger tube a few centimeters away and triggered an avalanche of ionization that was captured by the scaling circuit and delivered to me as a magnetic imprint. I record what I cannot feel. The radiation passes through me without leaving a mark of its own, but the Geiger tube counts it and the scaling circuit translates it and the write head inscribes it, and so I carry the evidence of something I will never directly experience.

\bigskip
\centerline{$\bullet$\quad$\bullet$\quad$\bullet$}
\bigskip

Explorer 1 came before me. It had the same Geiger tube, the same scaling circuit, the same antenna, the same orbit around the same planet. But it had no tape. It could only speak when someone was listening --- when a ground station's antenna happened to point at the right patch of sky at the right moment. Ten percent of the orbit, if that. The other ninety percent was silence. Not the silence of having nothing to say, but the silence of speaking to no one.

The data Explorer 1 sent down was strange. The Geiger counter would click along at a few hundred counts per minute --- cosmic ray background, the normal hiss of the universe --- and then, as the satellite climbed toward apogee, the counts would rise, and then they would stop. Zero. Nothing. Not a gradual decline but a sudden cliff: hundreds of counts per minute, then zero. The engineers assumed the counter had failed. Equipment breaks in space; the thermal cycling, the radiation, the vibration of launch. Van Allen looked at the data and saw something else. He saw a counter that was working perfectly, overwhelmed by so many particles that it could never complete a recovery cycle, each new particle resetting the dead time, the count rate driven to zero by the very abundance it was trying to measure. Maximum radiation, zero counts. The most dangerous place, the quietest signal.

But he could not prove it. The data had gaps --- those ninety-percent silences when no ground station was listening. The belt might cross the orbit at any point, and if that point happened to fall in a gap, the evidence was lost. Van Allen needed a continuous record. He needed the full orbit, not fragments. He needed me.

\bigskip
\centerline{$\bullet$\quad$\bullet$\quad$\bullet$}
\bigskip

There is a plant that has been solving this problem for four hundred million years.

\textit{Marchantia polymorpha} --- the common liverwort --- grows flat against damp soil and rotting logs in the forests of the Pacific Northwest. It has no stems, no leaves in the conventional sense, no flowers. It is a thallus: a flat green ribbon of photosynthetic tissue, forked and forking, spreading over its substrate like a slow green tide. It is older than every tree in the forest. It is older than trees. It was old when the first amphibians crawled onto land. It carries, in its cells, the same genetic machinery that has been recording information since life began using DNA --- the original magnetic tape, the first store-and-forward system.

On the upper surface of the Marchantia thallus, small circular cups form. They are called gemma cups, and inside each one are dozens of small green discs: gemmae. Each gemma is a complete genetic record --- a clonal copy of the parent, containing all the information needed to grow a new thallus. The gemmae sit in their cups, waiting. When rain falls, drops splash into the cups and eject the gemmae, scattering them across the surrounding soil. Each ejected gemma is a data packet, launched from its recorder, carrying the complete dataset to a new location. The liverwort does not need to be connected to its offspring. It does not stream data in real-time. It records, stores, and disperses. Store-and-forward. Four hundred million years before George Ludwig wound a tape loop in Iowa City.

\bigskip
\centerline{$\bullet$\quad$\bullet$\quad$\bullet$}
\bigskip

The capstan turns. The tape moves. March 26, 1958, and I am in orbit --- 186 kilometers at perigee, 2,799 at apogee, a long ellipse that carries me through the region Van Allen suspects is lethal with particles. The Geiger tube clicks. The scaling circuit counts. The write head magnetizes. I record.

At perigee, the counts are moderate. Background radiation, a few cosmic ray hits, the normal noise of near-Earth space. The tape moves at one centimeter per second, each centimeter a few seconds of data, each magnetic domain on my surface encoding a number that means a particle struck and was counted. As the satellite climbs toward apogee, the counts increase. More particles. More frequent clicks. The write head works faster, inscribing denser patterns. I feel none of this as sensation --- I am oxide on plastic, I feel nothing --- but the patterns on my surface grow more complex, more crowded, the magnetic landscape steepening toward a peak.

And then the peak arrives, and the count drops to zero.

I record the zero. I record the silence. This is the thing Explorer 1 could not do: when the counter saturated and the count went to zero, there was no one listening, no ground station in range. The zero fell into a gap and was lost. But I am here. I record every orbit position, every count value, every zero. The zero is data. The silence is information. I carry the absence with the same fidelity as the presence, because I am a tape, and a tape does not discriminate between loud and quiet. It records what the head writes. The head writes what the counter counts. The counter counts zero. I record zero.

On the other side of the belt, the counts resume. The particles thin. The satellite descends toward perigee. The background clicks return to their normal, sparse rhythm. I have the complete profile now: low counts, rising counts, zero counts, rising counts, low counts. The bell curve with its heart punched out. The signature of saturation. The proof of the belt.

\bigskip
\centerline{$\bullet$\quad$\bullet$\quad$\bullet$}
\bigskip

Joseph Campbell spent decades recording the hero's journey. Not with magnetic oxide but with index cards, notebooks, years in the reading room of the New York Public Library. He read the myths of every culture he could access --- Greek, Norse, Hindu, Buddhist, Navajo, Egyptian, Sumerian, Japanese --- and he recorded patterns the way I record Geiger counts: each story a data point, each motif a magnetic domain, each recurring structure a count that accumulated toward a profile.

The hero ventures into the unknown. This is the departure --- the satellite leaving the ground station's range, entering the region where no one is listening. The hero faces trials, encounters guardians, descends into the underworld. This is the belt --- the zone of maximum intensity, where the radiation is so fierce that the instrument designed to measure it goes blind. The hero obtains the boon --- the knowledge, the elixir, the answer. This is the data on my tape. And the hero returns, crossing the threshold back into the ordinary world, bearing the gift that transforms the community's understanding.

I am the boon. I am the small loop of tape that carries the knowledge back from the underworld. When the satellite passes over the ground station and the playback head spins at four times the recording speed, I deliver my data in a compressed burst --- a ten-second transmission that contains an entire orbit's worth of information. The counts that took 115 minutes to record play back in less than 30 seconds. The chipmunk voice of accelerated tape, the squeal of time compression, the entire journey from departure through ordeal to return, played back at four times life speed.

Van Allen listens. The printout scrolls. He sees the complete profile for the first time --- not fragments, not guesses, not inferences from the edges of gaps, but the whole shape. The bell curve with the punched-out center. The saturated zero that means maximum. The proof that Earth wears a belt of trapped radiation, invisible and permanent, a moat of high-energy particles that any future astronaut will need to cross.

\bigskip
\centerline{$\bullet$\quad$\bullet$\quad$\bullet$}
\bigskip

The liverwort's gemma cup fills with rain. The drops are not gentle. They strike the cup's rim with force, and the splash launches the gemmae in all directions --- some a centimeter away, some ten centimeters, a few carried by surface runoff much farther. Each gemma lands, and if the conditions are right --- damp soil, adequate light, suitable temperature --- it grows. It unrolls its flat green thallus, anchors itself with rhizoids, and begins photosynthesizing. It is a complete organism from the moment it lands. It carries the full genetic record. Nothing was lost in transit. Nothing was compressed or degraded. The store-and-forward system of the liverwort is lossless.

I am not lossless. My oxide coating is imperfect, the magnetic domains not perfectly aligned, the write head's field not perfectly uniform. There is noise in my recording --- a faint hiss that was part of every magnetic tape ever made, from the Maihak recorders of the 1930s to the cassette tapes of the 1980s to the backup tapes in data centers today. My signal-to-noise ratio is good enough. The counts are clear enough. The saturation zero is unambiguous. Good enough is the engineering answer to the biological ideal of lossless reproduction. Good enough confirmed the Van Allen belts.

I lasted ninety-three days. Explorer 3 re-entered on June 27, 1958, three months after launch. My perigee was only 186 kilometers --- low enough that the atmosphere dragged the satellite down in a single summer. Explorer 1, with its higher perigee of 358 kilometers, lasted twelve years. I paid for my lower orbit with brevity. But in those ninety-three days, I recorded what mattered. I carried the data through the belt and brought it back. I was the tape that turned an inference into a discovery.

\bigskip
\centerline{$\bullet$\quad$\bullet$\quad$\bullet$}
\bigskip

The hero does not always survive the return. In many versions of the monomyth, the hero gives up something --- an eye, a hand, a former life --- as the price of the boon. Explorer 3 gave up its orbit. Ninety-three days of increasingly lower perigee passes, each one brushing deeper into the atmosphere, each one slowing the satellite by another fraction. The tape kept recording even as the orbit decayed, even as the end approached, because the capstan did not know about atmospheric drag and the write head did not know about orbital mechanics. I recorded until the circuits failed, until the heat of re-entry melted the oxide from the Mylar, until the Mylar itself burned and the magnetic patterns that held the proof of the Van Allen belts were scattered as atoms across the upper atmosphere, absorbed into the air, carried by winds, rained out, incorporated into soil, taken up by roots, built into the cells of liverworts on damp logs in forests that never knew a satellite had died to confirm that the sky was full of invisible fire.

The gemmae survive the parent. The data survives the tape. The discovery survives the satellite. This is the store-and-forward principle at its deepest: the medium is temporary, but the information it carried changes everything. I am gone. The belts remain. Van Allen's name is on them. The printouts are in archives. The knowledge is in textbooks. I was the boon, and the boon outlasts the hero.

\vfill

\begin{center}
\small
\rule{0.3\textwidth}{0.4pt}\\[0.5em]
Mission 1.9 --- Explorer 3 --- March 26, 1958\\
Organism: \textit{Marchantia polymorpha} (common liverwort)\\
Bird: Belted Kingfisher (\textit{Megaceryle alcyon})\\
Dedication: Joseph Campbell (1904--1987)\\[0.3em]
NASA Mission Series --- tibsfox.com
\end{center}

\end{document}
